\documentclass[12pt]{article}
\usepackage[utf8]{inputenc}
\usepackage[T1]{fontenc}
\usepackage{hyperref}
\title{Andalusian Arabic Language Resources for K-12 Curriculum Development}
\author{Compiled by Claude Sonnet 5.5 (claude-sonnet-5-5) for Andr\'as Kornai}
\date{}
\begin{document}
\maketitle

\section{Introduction}
Andalusian Arabic (ISO 639-3: xaa) was the Arabic vernacular of al-Andalus, in what is now Spain and Portugal. It has no native speakers today. What survives is limited: mainly zajal poetry and glosses and lexical material. I found no school curriculum or revitalization programme in these searches, and no digital text corpus.

\section{Primary Educational Resources}
None found.

\section{Linguistic and Reference Materials}
\begin{itemize}
\item \emph{A Descriptive and Comparative Grammar of Andalusi Arabic} (Brill, 2013), from the Institute of Islamic Studies of the University of Zaragoza. I did not see a publisher URL. A review in the Journal of the American Oriental Society is at \url{https://lockwoodonlinejournals.com/index.php/jaos/article/view/1662}, and an Austrian Academy of Sciences library record is at \url{https://www.oeaw.ac.at/resources/Record/993582399004498}.
\item J. T. Monroe, article in \emph{Oral Tradition} arguing that the zajal is written in the Andalusi vernacular: \url{https://journal.oraltradition.org/wp-content/uploads/files/articles/4i-ii/4_Monroe.pdf}.
\item Overview of zajal poetry: \url{https://en.wikipedia.org/wiki/Zajal}.
\end{itemize}

\section{Religious and Community Resources}
None found.

\section{Digital and Audio Resources}
No audio, video or digital corpus was found. The Toponomasticon Hispaniae project cites the Zaragoza grammar for place-name etymology: \url{https://toponhisp.org/eu/node/5049}.

\section{Curriculum Development Context}
A heritage curriculum would be built from the Zaragoza grammar, zajal texts such as those of Ibn Quzman (c. 1078--1160), and Arabic-derived place names and loanwords in Iberian languages. Since there are no speakers, it would be a reconstruction exercise. It should be developed with specialists in Andalusi Arabic.

\end{document}